\scanpage{176}

\diagram{\begin{tikzcd}[ampersand replacement=\&,column sep=small]
\deleted{\unsure{\text{on}|\longrightarrow}}\ \mathrm{Prépic}'(X/S)\arrow[r]\&\mathrm{Prépic}'(X'/S')\arrow[r,shift left=.65ex]\arrow[r,shift right=.65ex]\&\mathrm{Prépic}'(X''/S'')
\end{tikzcd}}

[où $S''=S'\times_S S'$] est exacte, \deleted{i.e. la suite d'ens $\text{\uuline{Prépic}}_{X/S}$}

\deleted{\unsure{soit}} En effet, le point (i) est encore une conséquence immédiate de

\uline{Corollaire 1} Sous les conditions de prop. 1, \inserted{$X$ séparé quasi-cpt sur $S$} corollaire 1, soit $S'\longrightarrow S$ une extension fid plate \inserted{et qu. cpte} de la base. Soit $\underline E$ \deleted{\unsure{f}}faisceau \deleted{inversible}\inserted{loc. libre} sur $X$.

Pour qu'il soit isomorphe à l'image inverse d'un faisceau loc libre \deleted{[\unsure{rappelons que}]} sur $S$, il \unsure{f et s|faut} que $\underline E'$ soit isomorphe à l'image inverse d'un faisceau loc libre sur $S$.

En effet, il faut exprimer que $f_*(\deleted{\unsure{F}}\underline E)$ est loc. libre et $\deleted{\unsure{g}}f^*f_*(\underline E)\longrightarrow\underline E$ un isom. Or $f_*(\underline E)'=f_*(\underline E')$ \inserted{(cf \unsure{Künneth} \unsure{plat|plate})} et on est ramené à des assertions immédiates de stabilité de propriétés par extensions fid. plates (\unsure{Exp|Chap.} \uline{IV})

% archive sheet number: 175
